\chapter{Encoder Eğitimi}
\label{ch:egitim}

Buraya kadarki bölümler sabit ağırlıklı bir checkpoint'i hızlandırmayı
anlattı. EdgeTAM indirildiği gibi alındı, dört TensorRT motoruna taşındı,
ölçüldü. Bu bölüm ters taraf: ağırlıkların kendisini termal ve havadan
veriyle eğitmek.

Hat üç aşamalı. \textbf{A} ve \textbf{B} kurulu ve kodu çalışır hâlde;
sıradaki iş onları koşturmak. \textbf{C} önerilen ek adım ve B ölçülmeden
başlamıyor.

Bölüm kasten görsel. Her şekil bir kurulumu anlatıyor, metin şekli
açıklıyor.

\section{Ne Eğitiliyor}
\label{sec:egitim-dort-parca}

EdgeTAM dört parçadan oluşuyor. İkisi tek bir kareyi maskeye çeviriyor.
Diğer ikisi hafıza yolu, yani modeli tracker yapan kısım.

\begin{figure}[H]
  \centering
  \begin{tikzpicture}[font=\elyazi\small]
    \eskizbasla[11]
    % --- üst sıra: tek kare, görüntüden maskeye ---
    \node[kgri,  text width=15mm] (im)   at (0,0)    {görüntü};
    \node[kmavi, text width=25mm] (enc)  at (3.0,0)
      {image encoder\\{\footnotesize RepViT + FPN}\\{\footnotesize \textbf{4,92 M}}};
    \node[ksari, text width=21mm] (feat) at (6.3,0)
      {feature map\\{\footnotesize $256{\times}32{\times}32$}};
    \node[kyesil, text width=20mm] (dec) at (9.2,0)  {mask\\decoder};
    \node[kgri,  text width=14mm] (mask) at (12.3,0) {maske};

    \draw[esok] (im)   -- (enc);
    \draw[esok] (enc)  -- (feat);
    \draw[esok] (feat) -- (dec);
    \draw[esok] (dec)  -- (mask);

    % prompt encoder: mask decoder'ın soluna kaydırıldı ki aşağıdaki
    % decoder -> memory encoder hattı içinden geçmesin
    \node[kyesil, text width=24mm] (pe) at (7.4,-2.1)
      {prompt encoder\\{\footnotesize box / point}};
    \draw[esok] (pe) -- (dec);

    % --- alt sıra: hafıza yolu, frozen ---
    \node[kdonuk, text width=25mm] (ma)   at (2.6,-4.1)
      {memory attention\\{\footnotesize 2,96 M}};
    \node[kdonuk, text width=22mm] (bank) at (6.3,-4.1) {memory bank};
    \node[kdonuk, text width=25mm] (me)   at (10.2,-4.1)
      {memory encoder\\{\footnotesize 1,62 M}};

    \draw[esokr=esgri] (me)   -- (bank);
    \draw[esokr=esgri] (bank) -- (ma);
    \draw[esokr=esgri] (ma.north) -- (feat.south);
    % decoder -> memory encoder: mask decoder'ın sağ alt kenarından dik
    % iniyor, hiçbir kutunun içinden geçmeden
    \draw[esokr=esgri] (10.2,-0.5) -- (me.north);

    % işaret
    \node[esisaret, fit=(enc)(feat)(dec)(pe)] (grup) {};
    \node[esvurgu, anchor=south, yshift=1mm] at (grup.north) {eğitilen kısım};
    \node[esnot, anchor=north, yshift=-2mm, text width=95mm] at (6.3,-4.9)
      {taralı kutular \textbf{frozen}: gradient almıyor, yalnızca kareler
       arasında çalışıyor};
  \end{tikzpicture}
  \caption{EdgeTAM'ın dört parçası. Üst sıra tek bir kareyi maskeye
    çeviriyor, alt sıra kareler arasında hafıza taşıyor. Bu bölümdeki
    eğitimin tamamı üst sırada geçiyor.}
  \label{fig:egitim-dort-parca}
\end{figure}

Hafıza yolu frozen kalıyor. Üç sebep var, herhangi biri tek başına yeterdi:

\begin{enumerate}
  \item \textbf{Recurrent bir döngünün yazma ucu.} Bu proje oradaki bir
        değişikliğin ne yaptığını zaten ölçtü: yalnızca memory encoder'ı
        nicemlemek ortalama IoU'yu 0,9397'ye düşürdü. Hasar tek karede
        kalmadı, klip boyunca birikti; çünkü onun çıktısı memory bank'a
        yazılıp sonraki yedi karede geri okunuyor
        (\diger{\cref{ch:deneyler}}{ana raporun 4. bölümü}).
  \item \textbf{ONNX yeniden yazımının zor olanı orası:} sabit memory
        slotları, döşenmiş RoPE tabloları, additive mask. O grafiklerin
        türetildiği ağırlıkları yeniden eğitmek, checkpoint ile motor
        arasında sessiz bir uyuşmazlık riski demek.
  \item \textbf{Piksel değil feature işliyor.} Termal ile RGB arasındaki
        fark bir encoder problemi, yani memory yolunun termale özel bir
        eğitime ihtiyacı yok. Encoder taşındıktan sonra ona \emph{yeniden}
        hizalanması gerekip gerekmediği ayrı bir soru; bir sonraki
        bölümün son kısmı onu ele alıyor.
\end{enumerate}

SAM~2 de aynı ayrımı yapıyor. Encoder'ı statik görüntüyle, hafıza yolunu
videoyla eğitiyor; uzun dizilere geçerken image encoder'ı tamamen
donduruyor~\cite{ravi2024sam2}.

\section{Üç Aşama}
\label{sec:egitim-uc-asama}

\begin{figure}[H]
  \centering
  \begin{tikzpicture}[font=\elyazi\small]
    \eskizbasla[42]
    % A
    \node[kmavi, text width=31mm] (a) at (0,0)
      {\textbf{A: pretraining}\\[0.5mm]{\footnotesize image encoder}};
    \node[ksari, text width=31mm] (ad) at (0,-2.0)
      {{\footnotesize hizalı RGB-T çifti}\\{\footnotesize \textbf{etiketsiz}}};
    \node[kgri, text width=31mm] (al) at (0,-3.7)
      {{\footnotesize cosine loss}};
    \draw[esok] (a) -- (ad);  \draw[esok] (ad) -- (al);

    % B
    \node[kmavi, text width=31mm] (b) at (4.7,0)
      {\textbf{B: fine-tune}\\[0.5mm]{\footnotesize image encoder
       + mask decoder}};
    \node[ksari, text width=31mm] (bd) at (4.7,-2.0)
      {{\footnotesize statik kare +}\\{\footnotesize \textbf{dense mask}}};
    \node[kgri, text width=31mm] (bl) at (4.7,-3.7)
      {{\footnotesize 20·focal + dice + IoU}};
    \draw[esok] (b) -- (bd);  \draw[esok] (bd) -- (bl);

    % C
    \node[kdonuk, text width=31mm] (c) at (9.4,0)
      {\textbf{C: video}\\[0.5mm]{\footnotesize memory attention
       + memory encoder}};
    \node[kdonuk, text width=31mm] (cd) at (9.4,-2.0)
      {{\footnotesize video + instance ID}};
    \node[kdonuk, text width=31mm] (cl) at (9.4,-3.7)
      {{\footnotesize + object-score BCE}};
    \draw[esokr=esgri] (c) -- (cd);  \draw[esokr=esgri] (cd) -- (cl);

    \draw[esok]  (a) -- (b);
    \draw[esink] (b) -- (c);

    \node[esiyi, anchor=south, yshift=2.5mm] at (a.north) {kurulu};
    \node[esiyi, anchor=south, yshift=2.5mm] at (b.north) {kurulu};
    \node[esnot, anchor=south, yshift=2.5mm] at (c.north) {önerilen ek adım};

    \node[esnot, anchor=north, yshift=-2.5mm, text width=62mm] at (2.35,-4.2)
      {sıradaki iş: bu ikisini koşturmak};
  \end{tikzpicture}
  \caption{Üç aşama, her biri kendi verisi ve kendi loss'uyla. Aşama A hiç
    etiket okumuyor, B dense mask istiyor, C instance ID'li video istiyor.
    Sıra keyfi değil: SAM~2'nin tarifi de önce statik görüntü, sonra video.}
  \label{fig:egitim-uc-asama}
\end{figure}

Her aşamada başka bir şey açılıyor, gerisi frozen kalıyor:

\begin{center}\footnotesize
\begin{tabular}{@{}l>{\raggedright\arraybackslash}p{95mm}@{}}
  \toprule
  \textbf{aşama} & \textbf{eğitime açılan} \\
  \midrule
  A & yalnızca image encoder \\
  B & önce mask decoder + prompt encoder, sonra bunlara image encoder da
      katılıyor \\
  C & memory attention + memory encoder; image encoder bu kez frozen \\
  \bottomrule
\end{tabular}
\end{center}

\subsection*{Aşama C neden var}

Memory attention'ın $W_Q$ ve $W_K$ matrisleri, EdgeTAM eğitilirken image
encoder'ın ürettiği feature dağılımına oturdu. Aşama A ve B o dağılımı
taşıyor. Mask decoder aşama B'de encoder'la \emph{birlikte} eğitildiği için
yeni dağılıma uyum sağlıyor. Memory yolu sağlamıyor: aşama B tek kare
üzerinde ve boş bir memory bank'la koşuyor, yani oraya hiç gradient
gitmiyor.

\begin{figure}[H]
  \centering
  \begin{tikzpicture}[font=\elyazi\small]
    \eskizbasla[137]
    \node[kmavi,  text width=23mm] (enc) at (1.6,2.6)
      {image encoder\\{\footnotesize A ve B'de hareket etti}};
    \node[ksari,  text width=20mm] (feat) at (5.0,2.6)
      {feature map\\{\footnotesize yeni uzay}};
    \node[kkirmizi, text width=25mm] (ma) at (8.6,2.6)
      {memory attention\\{\footnotesize eski uzaya oturdu}};
    \node[kyesil, text width=20mm] (dec) at (12.2,2.6)
      {mask decoder\\{\footnotesize B'de uyum sağladı}};

    \node[kkirmizi, text width=25mm] (me) at (5.0,0.3)
      {memory encoder\\{\footnotesize eski uzaya oturdu}};
    \node[kgri, text width=21mm] (bank) at (8.6,0.3) {memory bank};

    \draw[esok] (enc) -- (feat);
    \draw[esok] (feat) -- node[esnot, anchor=south, yshift=0.5mm] {Q} (ma);
    \draw[esok] (ma) -- (dec);
    \draw[esok] (feat) -- node[esnot, anchor=west, xshift=1.5mm] {+ maske} (me);
    \draw[esok] (me) -- (bank);
    \draw[esok] (bank) -- node[esnot, anchor=west, xshift=1.5mm] {K, V} (ma);

    \node[esnot, anchor=north, yshift=-4mm, text width=120mm] at (6.8,-0.3)
      {kırmızı kutular A ve B boyunca hiç gradient almıyor. Aşama C tam olarak
       onları açıyor, encoder'ı bu kez donduruyor.};
  \end{tikzpicture}
  \caption{Aşama A ve B bittiğinde encoder'ın yeni feature uzayını görmemiş
    tek yer memory yolu. Aşama C'nin işi yeni bir yetenek öğretmek değil, o
    iki kutuya yeni uzayı okumayı yeniden öğretmek.}
  \label{fig:egitim-hiza}
\end{figure}

Bu beklenen bir çöküş değil. Attention'ın iki tarafı da aynı encoder'dan
iniyor: Q geçerli karenin feature'ları, K ve V ise memory encoder'ın
\emph{aynı} feature'lardan yazdığı memory. Kayma tek tarafa değil ikisine
birden uygulanıyor. Ama $W_Q$ ve $W_K$ sabit ve kaymanın benzerlik yapısını
koruduğunu garanti eden bir şey yok. Yani bozulmanın büyüklüğü ölçülene
kadar bilinmiyor.

Ölçen sayı da henüz elimizde değil. Elimizdeki metrik tek prompt'lu tek kare
skorluyor ve memory hiç devreye girmiyor, yani bu hasarı tanım gereği
göremiyor. Görebilen tek metrik J\&F, o da aşama C'nin kurulumunu istiyor.

\section{Aşama A}
\label{sec:egitim-asama-a}

Aşama A'nın tek varlık sebebi şu: hizalı RGB-T çiftleri hiçbir etiket
istemiyor, ve elimizde onlardan çok var. Bir RGB foundation model teacher
oluyor, EdgeTAM'ın termal encoder'ı student. Teacher çiftin RGB yarısına
bakıyor, student termal yarısına~\cite{anythermal2026}.

\begin{figure}[H]
  \centering
  \begin{tikzpicture}[font=\elyazi\small]
    \eskizbasla[7]
    \node[kgri, text width=19mm] (rgb)  at (0,1.0)  {RGB kare};
    \node[kgri, text width=19mm] (ter)  at (0,-1.5) {termal kare};
    \node[esnot, text width=22mm] at (0,-0.25) {\emph{hizalı çift}};

    \node[kdonuk, text width=29mm] (t) at (3.6,1.0)
      {teacher: DINOv3\\{\footnotesize ViT-B/16, \textbf{frozen}}};
    \node[kmavi, text width=29mm] (s) at (3.6,-1.5)
      {student: EdgeTAM\\{\footnotesize image encoder}};

    \draw[esokr=esgri] (rgb) -- (t);
    \draw[esok] (ter) -- (s);

    \node[ksari, text width=23mm] (tf) at (7.3,1.0)
      {teacher feature\\{\footnotesize $768{\times}32{\times}32$}};
    \node[ksari, text width=23mm] (sf) at (7.3,-1.5)
      {student feature\\{\footnotesize $256{\times}32{\times}32$}};
    \draw[esokr=esgri] (t) -- (tf);
    \draw[esok] (s) -- (sf);

    \node[kmor, text width=21mm] (proj) at (10.7,-1.5)
      {$1{\times}1$ projection\\{\footnotesize $256 \to 768$}};
    \draw[esok] (sf) -- (proj);

    \node[kyesil, text width=23mm] (loss) at (10.7,1.0)
      {cosine loss\\{\footnotesize pozisyon pozisyon}};
    \draw[esokr=esgri] (tf) -- (loss);
    \draw[esok] (proj) -- (loss);

    \node[esvurgu, anchor=north, yshift=-1.5mm, text width=34mm]
      at (proj.south) {aşama bitince atılıyor};
    \node[esiyi, anchor=north, yshift=-3mm, text width=38mm]
      at (s.south) {hiçbir etiket okunmuyor};
  \end{tikzpicture}
  \caption{Aşama A. Teacher frozen, student EdgeTAM'ın kendi image
    encoder'ı. Aradaki $1{\times}1$ projection yalnızca iki kanal sayısını
    hizalıyor ve aşama bitince atılıyor. Çıkan dosya sıradan bir EdgeTAM
    checkpoint'i: aynı anahtarlar, aynı loader, aynı ONNX grafiği.}
  \label{fig:egitim-asama-a}
\end{figure}

\subsection*{Neden MAE değil}

Akla ilk gelen pretraining yöntemi MAE, çünkü SAM~2'nin Hiera backbone'u
öyle eğitilmiş. Ama transfer olmuyor. MAE token maskeliyor ve ViT için
tasarlandı; EdgeTAM'ın backbone'u convolutional
RepViT~\cite{wang2023repvit}, ortada maskelenecek bir token grid'i yok.

Feature distillation'ın böyle bir şartı yok. Teacher ile student'ın
mimarileri ilgisiz olabilir. Hizalanması gereken tek şey bir feature map,
onu da $1{\times}1$ projection hallediyor.

\subsection*{Teacher seçimi}

Önce bir yanlış anlamayı kapatalım: \textbf{teacher hiçbir zaman EdgeTAM'ın
içine takılmıyor.} Ayrı bir model olarak koşuyor, bir feature map üretiyor,
ve karşılaştırılan tek şey sayılar. Mimarilerin ilgisiz olması bir sorun
değil, distilasyonun tanımı. Bu yüzden Orin'de asla koşamayacak kadar büyük
bir model bile burada sorunsuz teacher olabiliyor: o çevrimdışı koşuyor,
dağıtılan şey EdgeTAM'ın kendi 4,92 M'lik encoder'ı.

Üç aday var. Aralarındaki en somut fark, teacher'ın feature map'inin
loss'tan önce resample edilip edilmediği.

\begin{figure}[H]
  \centering
  \begin{tikzpicture}[font=\elyazi\small]
    \eskizbasla[19]
    \node[kmor, text width=34mm] (d3) at (0,0)
      {\textbf{DINOv3} ViT-B/16\\{\footnotesize $512/16 = 32{\times}32$}\\
       {\footnotesize \emph{kod varsayılanı}}};
    \node[kgri, text width=34mm] (d2) at (5.3,0)
      {DINOv2 base\\{\footnotesize $518/14 = 37{\times}37$}\\
       {\footnotesize açık, hesap gerekmez}};
    \node[kgri, text width=34mm] (s2) at (10.6,0)
      {SAM 2.1 Hiera-L\\{\footnotesize $1024/16 = 64{\times}64$}\\
       {\footnotesize SAM 2'nin encoder'ı}};

    \node[ksari, text width=40mm] (g) at (5.3,-4.0)
      {student grid\\{\footnotesize $512/16 = \mathbf{32\times32}$}};

    \draw[esokr=esyesil] (0,-1.15) -- (0,-4.0) -- (g.west);
    \draw[esink] (d2) -- (g);
    \draw[esokr=esgri] (10.6,-1.15) -- (10.6,-4.0) -- (g.east);

    \node[esiyi, anchor=west, xshift=1.5mm, text width=32mm] at (0,-2.4)
      {\textbf{birebir}\\resample yok};
    % iki "interpolasyon" etiketi: biri kendi çizgisinin sağında, öteki
    % solunda -- aynı yükseklikte oldukları için aksi hâlde üst üste binerler
    \node[esnot, anchor=east, xshift=-1.5mm, text width=24mm] at (10.6,-2.4)
      {interpolasyon};
    \node[esnot, anchor=west, xshift=1.5mm, text width=24mm] at (5.3,-2.4)
      {interpolasyon};
  \end{tikzpicture}
  \caption{Teacher adaylarının grid'i, student'ınkiyle karşılaştırmalı.
    DINOv3'ün 16 piksellik patch'i 512 girdide tam olarak student grid'ini
    veriyor, yani teacher'ın feature map'i loss'tan önce hiç resample
    edilmiyor. Varsayılanın DINOv3 olmasının asıl sebebi ise grid değil:
    dense feature kalitesi o sürümün konusu~\cite{dinov32025}.}
  \label{fig:egitim-ogretmen}
\end{figure}

Üçü de kodda hazır, seçim tek bir string, ve aralarında ölçülmüş bir
karşılaştırma yok. Asıl gerilim DINOv3 ile SAM 2.1 arasında:

\begin{itemize}
  \item \textbf{DINOv3 lehine:} aradığımız şey anlam. Termalde sınır zaten
        kolay, soğuk arka plandaki sıcak nesne belirgin bir kenar; eksik olan
        o kenarın \emph{ne} olduğu. SAM 2'nin feature'ları ise tasarım gereği
        sınıftan bağımsız, her şeyi segmentliyor ama ne olduğunu bilmiyor.
  \item \textbf{SAM 2.1 lehine:} üç şey birden. Tensörler birebir uyuyor,
        256 kanal ve stride 16, yani projection bir temsil değişikliği değil.
        Encoder'ı SAM 2'nin uzayına \emph{doğru} itiyor, ondan uzağa değil,
        yani bir önceki bölümdeki hiza riskini büyütmek yerine küçültüyor. Ve
        en önemlisi: EdgeTAM \emph{zaten} SAM 2'nin distilasyonu, dolayısıyla
        ondan ``termal girdiden SAM 2'nin RGB feature'larını üret'' demek,
        checkpoint'in doğduğu hedefin aynısını istemek oluyor, üstüne yalnızca
        modalite boşluğu eklenmiş hâlde. Bedeli hesap: Hiera-Large 1024'te
        koşuyor, ki EdgeTAM'ın var olma sebebi tam olarak ondan kaçmak.
\end{itemize}

\textbf{İkisi birden kullanılmıyor,} ve sebebi teknik bir engel değil.
Çoklu-teacher yayınlanmış bir desen, ama burada iki loss iki ayrı aşamada
ve verileri iki mertebe farklı; ortak bir batch aşama A'nın çoğunu çöpe
atardı. Aynı fikrin ucuz hâli zaten aşağıda: çapa terimi, yani feature
distillation'ı ikinci bir teacher olarak değil regularizer olarak
kullanmak.

\textbf{Önce SAM 2.1 koşulacak.} Sebep DINOv3'ün zayıf olması değil,
sıralamanın riski: SAM 2.1 varyantı aşama C hiç yapılmasa bile ayakta kalan
varyant, çünkü encoder'ı memory yolunun zaten beklediği uzayda tutuyor.
DINOv3 varyantının üst sınırı daha yüksek ama hiza riskini o taşıyor, yani
ikinci koşu olması doğru. Aradaki tek fark bir string; veri, seed, schedule
ve loss aynı kalıyor, dolayısıyla çıkan fark teacher'ın kendisine ait
oluyor.

\textbf{Hangisi olursa olsun, hiçbir şey yapmamayı yenmek zorunda.}
Baseline: stok EdgeTAM checkpoint'inden başlayan aşama B. Pretraining bu
baseline'ı geçmiyorsa harcanmış GPU saatidir. Karşılaştırılabilir olan tek
sayı test split'indeki instance IoU'su; cosine loss teacher'lar arasında
karşılaştırılamaz.

Aşama B, aşama A'nın öğrendiğinden geri kaymasın diye bir \textbf{çapa
terimi} taşıyor: encoder'ın feature'ları aşama B başındaki hâlinden ne kadar
uzaklaştıysa o kadar ceza. Bu çapa aşama A'yı koruyor, memory yolunun
hizasını değil.

\subsection*{4,92 M parametre iki modaliteye yeter mi}

Encoder küçük, ve şu an yalnızca termal öğreniyor. RGB yarısı da eğitime
katılırsa aynı ağırlıklar iki modaliteyi birden taşımak zorunda kalır. İki
şey bunu göründüğünden ucuz yapıyor: encoder sıfırdan öğrenmiyor, zaten
SAM~2'nin distilasyonu; ve iki modalite ilgisiz alanlar değil, aynı sahneye
aynı açıdan bakıyorlar. Yine de kapasite gerçek bir sınır ve ölçülmedi. En
ucuz test: iki yarıyı da eğitip termal IoU'yu yalnızca termalle eğitilmiş
koşuya karşı okumak. Düşerse sınır kapasitedir; düşmezse RGB bedava
gelmiştir.

\section{Aşama B}
\label{sec:egitim-semantik}

Havadan dense mask'li setlerin neredeyse tamamı \textbf{semantic
segmentation} veriyor: piksel $\rightarrow$ sınıf, yani ``bütün arabalar
araba''. Tracker'ın istediği bunun tersi: ``şu araba, yanındaki değil''.

Semantic bir target'la eğitmek nötr değil, zararlı. Yan yana iki arabanın
feature'larını birbirine yaklaştırır; oysa tracking için ayrışmaları
gerekiyor.

Çözüm veri setini değiştirmek değil, aynı veri üzerinde target'ı
değiştirmek. SAM~2 de zaten semantic map görmüyor: görüntü başına tek tek
mask örnekleyip her birini ayrı prompt ediyor.

\begin{figure}[H]
  \centering
  \begin{tikzpicture}[font=\elyazi\small]
    \eskizbasla[31]
    % --- 1. semantic map ---
    \begin{scope}[shift={(0,0)}]
      \foreach \x/\y/\c/\t in {%
        0/2/esgrid/yol, 1/2/esgrid/yol, 2/2/esgrid/yol,
        0/1/esgrid/yol, 1/1/eskirmizid/ARB, 2/1/eskirmizid/ARB,
        0/0/esmord/bina, 1/0/eskirmizid/ARB, 2/0/eskirmizid/ARB}
        {\draw[cizik, draw=escizgi, fill=\c] (\x*0.9,\y*0.68) rectangle ++(0.9,0.68);
         \node[font=\elyazi\tiny, text=esmurekkep] at (\x*0.9+0.45,\y*0.68+0.34) {\t};}
      \node[esbaslik, anchor=north, yshift=-2mm] at (1.35,0) {semantic map};
    \end{scope}

    % Ok etiketleri okun *altında*: ızgaralar y=0..2,04 arasını kaplıyor, o
    % yüzden okun üstünde kalan bir etiket kaçınılmaz olarak onlara biner.
    \draw[esok] (2.9,1.35) -- (5.3,1.35);
    \node[esnot, anchor=north, yshift=-1.5mm, text width=24mm] at (4.1,1.2)
      {stuff sınıflarını at};

    % --- 2. filtrelenmiş ---
    \begin{scope}[shift={(5.4,0)}]
      \foreach \x/\y/\c/\t in {%
        0/2/white/{$\cdot$}, 1/2/white/{$\cdot$}, 2/2/white/{$\cdot$},
        0/1/white/{$\cdot$}, 1/1/eskirmizid/ARB, 2/1/eskirmizid/ARB,
        0/0/white/{$\cdot$}, 1/0/eskirmizid/ARB, 2/0/eskirmizid/ARB}
        {\draw[cizik, draw=escizgi, fill=\c] (\x*0.9,\y*0.68) rectangle ++(0.9,0.68);
         \node[font=\elyazi\tiny, text=esmurekkep] at (\x*0.9+0.45,\y*0.68+0.34) {\t};}
      \node[esbaslik, anchor=north, yshift=-2mm, text width=32mm] at (1.35,0)
        {yalnızca thing sınıfları};
    \end{scope}

    \draw[esok] (8.3,1.35) -- (10.7,1.35);
    \node[esnot, anchor=north, yshift=-1.5mm, text width=26mm] at (9.5,1.2)
      {sınıf sınıf connected component};

    % --- 3. iki instance ---
    \begin{scope}[shift={(10.9,0)}]
      \node[kyesil, text width=8mm, minimum height=10mm] at (0.6,1.35) {\textbf{1}};
      \node[kyesil, text width=8mm, minimum height=10mm] at (2.2,1.35) {\textbf{2}};
      \node[esbaslik, anchor=north, yshift=-2mm, text width=32mm] at (1.4,0.6)
        {iki instance, iki prompt, iki target};
    \end{scope}
  \end{tikzpicture}
  \caption{Semantic map'ten instance target'ına. Thing sınıfları Kust4K'nın
    sekizinden dördü: motosiklet, araba, kamyon, insan. Yol, bina ve ağaç
    stuff, yani takip hedefi değil.}
  \label{fig:egitim-semantik-ornek}
\end{figure}

\subsection*{Bunun kendi riski: kaynaşma}

Connected component geometriyle karar veriyor. Tampon tampona park etmiş
iki araba tek bir piksel kenarı paylaşıyorsa tek component olur, ve hiçbir
eşik bunu ayıramaz.

\begin{figure}[H]
  \centering
  \begin{tikzpicture}[font=\elyazi\small]
    \eskizbasla[57]
    % sol: ince köprü -> watershed kurtarır
    \begin{scope}[shift={(0,0)}]
      \draw[cizik, draw=eskirmizi, fill=eskirmizid] (0,0) rectangle (1.5,1.0);
      \draw[cizik, draw=eskirmizi, fill=eskirmizid] (1.85,0) rectangle (3.35,1.0);
      \draw[cizik, draw=eskirmizi, fill=eskirmizid] (1.45,0.42) rectangle (1.9,0.58);
      \node[esbaslik, anchor=north, yshift=-2.5mm, text width=44mm] at (1.67,0)
        {\textbf{ince köprü}\\{\footnotesize distance transform'da iki tepe,
         arada vadi}};
      \node[esiyi, anchor=north, yshift=-2.5mm, text width=44mm] at (1.67,-1.75)
        {watershed ayırıyor};
    \end{scope}
    % sağ: tam kenar -> hiçbir şey kurtaramaz
    \begin{scope}[shift={(6.6,0)}]
      \draw[cizik, draw=eskirmizi, fill=eskirmizid] (0,0) rectangle (1.65,1.0);
      \draw[cizik, draw=eskirmizi, fill=eskirmizid] (1.65,0) rectangle (3.3,1.0);
      \node[esbaslik, anchor=north, yshift=-2.5mm, text width=46mm] at (1.65,0)
        {\textbf{tam kenar boyunca bitişik}\\{\footnotesize daha büyük bir
         dikdörtgene döşenir, vadi yok}};
      \node[esvurgu, anchor=north, yshift=-2.5mm, text width=48mm] at (1.65,-1.75)
        {mask geometrisiyle çalışan \emph{hiçbir} yöntem ayıramaz};
    \end{scope}
  \end{tikzpicture}
  \caption{Watershed'in sınırı, ve bu sınır denenmeden önce bilinmesi
    gereken şey.}
  \label{fig:egitim-kaynasma}
\end{figure}

Bu riskin ölçüsü var, varsayımı yok. Bir component'in kendi kutusunu ne
kadar doldurduğu kaynaşmanın sinyali: dörtte birini dolduran bir component
genelde ince bir piksel köprüsüyle birleşmiş iki nesnedir. Notebook bu
oranı \textbf{tek bir GPU saati harcanmadan önce} sayıyor, ve o sayı bütün
B aşamasının git/gitme kararı.

Kötü çıkarsa dört çıkış yolu var, maliyet sırasıyla: \textbf{(1)} watershed
modunu aç, yalnızca köprü vakasını kurtarır ve bedavadır; \textbf{(2)}
zaten instance mask veren bir set kullan, ki VTUAV VIS'te zaten var;
\textbf{(3)} doldurma eşiğini sertleştir, daha az ama daha temiz instance;
\textbf{(4)} \textbf{SAM'a sor.} Sonuncusu en iyisi çünkü karar
\emph{görünüşe} göre veriliyor, mask geometrisine göre değil: SAM~3'e metin
prompt'u veriliyor (``car'') ve o, bitişik iki arabayı bile ayrı ayrı
maskeliyor~\cite{sam32025}. Eğitim döngüsünün dışında, çevrimdışı ve bir
kez koşuyor.

\subsection*{Bir batch nasıl kuruluyor}

Encoder prompt'a bağlı değil, yani bir görüntüyü kaç kez prompt ettiğiniz
onun maliyetini değiştirmiyor. O yüzden görüntü bir kez encode ediliyor ve
pencere içindeki her instance aynı feature'lara karşı prompt ediliyor.

\begin{figure}[H]
  \centering
  \begin{tikzpicture}[font=\elyazi\small]
    \eskizbasla[71]
    \draw[cizik, draw=escizgi, fill=white] (0,-1.4) rectangle (3.2,1.4);
    \node[esbaslik, anchor=north, yshift=-2mm, text width=40mm] at (1.6,-1.4)
      {pencere\\{\footnotesize $512\times512$ native piksel}};
    \draw[cizik, draw=eskirmizi, fill=eskirmizid] (0.45,0.55) rectangle (1.35,0.95);
    \node[font=\elyazi\tiny] at (0.9,0.75) {A};
    \draw[cizik, draw=eskirmizi, fill=eskirmizid] (1.85,-0.35) rectangle (2.75,0.05);
    \node[font=\elyazi\tiny] at (2.3,-0.15) {B};
    \draw[cizik, draw=eskirmizi, fill=eskirmizid] (0.6,-0.95) rectangle (1.1,-0.6);
    \node[font=\elyazi\tiny] at (0.85,-0.78) {C};

    \node[kmavi, text width=24mm] (enc) at (5.2,0)
      {image encoder\\{\footnotesize \textbf{tek geçiş}}};
    \draw[esok] (3.35,0) -- (enc);

    \node[ksari, text width=19mm] (f) at (7.9,0) {feature map};
    \draw[esok] (enc) -- (f);

    \node[kyesil, text width=30mm] (d1) at (11.1,1.25)
      {decoder $\leftarrow$ box A};
    \node[kyesil, text width=30mm] (d2) at (11.1,0)
      {decoder $\leftarrow$ box B};
    \node[kyesil, text width=30mm] (d3) at (11.1,-1.25)
      {decoder $\leftarrow$ box C};
    \draw[esok] (f) -- (d1);
    \draw[esok] (f) -- (d2);
    \draw[esok] (f) -- (d3);

    \node[esiyi, anchor=north, yshift=-3mm, text width=76mm] at (8.0,-1.7)
      {batch artık aynı sahneyi paylaşan birden çok nesne taşıyor. Loss'un
       yapması gereken ayrım tam olarak bu.};
  \end{tikzpicture}
  \caption{Tek encode, instance başına bir prompt. Hem daha ucuz hem daha
    iyi sinyal: aynı sahnedeki iki aracı ayırmak, loss'un öğrenmesi gereken
    şeyin kendisi.}
  \label{fig:egitim-batch}
\end{figure}

İki ayrıntı, sessiz hata üreten türden. \textbf{Split kare üzerinden
yapılıyor, pencere üzerinden değil:} bir görüntünün iki penceresi
piksellerinin çoğunu paylaşır, pencere seviyesinde bölmek neredeyse aynı
görüntüleri çizginin iki yanına koyar ve held-out sayı anlamını kaybeder.
\textbf{Kenardan kesilen instance kırpılmıyor, düşürülüyor:} piksellerinin
sahip olmadığı bir genişliği iddia eden box, hedefsizlikten daha kötü bir
target'tır.

\section{Veri}
\label{sec:egitim-veri}

Üç set kullanılıyor ve üçü de aynı işi yapmıyor. Aşağıdaki tablo bütün
hikâye.

\begin{figure}[H]
  \centering
  \begin{tikzpicture}[font=\elyazi\small]
    \eskizbasla[151]
    % --- sütun başlıkları ---
    \node[esbaslik, text width=34mm] at (4.5,2.35)  {\textbf{A} \\{\footnotesize ön-eğitim}};
    \node[esbaslik, text width=34mm] at (8.4,2.35)  {\textbf{B} \\{\footnotesize fine-tune}};
    \node[esbaslik, text width=34mm] at (12.3,2.35) {\textbf{C} \\{\footnotesize video, önerilen}};

    % --- VTUAV ---
    \node[esbaslik, anchor=east, text width=22mm] at (2.4,1.0)
      {\textbf{VTUAV}\\{\footnotesize $1920{\times}1080$}};
    \node[kmavi,  text width=32mm] at (4.5,1.0)
      {{\footnotesize \textbf{$\approx$1 700 000 çift}}\\{\footnotesize etiket okunmuyor}};
    \node[kyesil, text width=32mm] at (8.4,1.0)
      {{\footnotesize VIS bölümü: 100 dizi}\\{\footnotesize \textbf{gerçek instance mask}}};
    \node[kmor,   text width=32mm] at (12.3,1.0)
      {{\footnotesize kutu etiketli 400 dizi}\\{\footnotesize SAM ile mask'e çevrilir}};

    % --- SegFly ---
    \node[esbaslik, anchor=east, text width=22mm] at (2.4,-0.55)
      {\textbf{SegFly}\\{\footnotesize $640{\times}512$}};
    \node[kmavi,  text width=32mm] at (4.5,-0.55)
      {{\footnotesize 15 007 çift}};
    \node[kyesil, text width=32mm] at (8.4,-0.55)
      {{\footnotesize semantic}\\{\footnotesize $\rightarrow$ component}};
    \node[esnot] at (12.3,-0.55) {video yok};

    % --- Kust4K ---
    \node[esbaslik, anchor=east, text width=22mm] at (2.4,-1.9)
      {\textbf{Kust4K}\\{\footnotesize $640{\times}512$}};
    \node[kmavi,  text width=32mm] at (4.5,-1.9)
      {{\footnotesize 4 024 çift}};
    \node[kyesil, text width=32mm] at (8.4,-1.9)
      {{\footnotesize semantic}\\{\footnotesize $\rightarrow$ component}};
    \node[esnot] at (12.3,-1.9) {video yok};
  \end{tikzpicture}
  \caption{Hangi set hangi aşamayı besliyor. Bir sette \emph{ne} olduğu
    değil, o aşamanın \emph{ne okuduğu} belirliyor: aşama A yalnızca hizalı
    çift sayıyor, aşama B dense mask istiyor, aşama C kimlikli video.}
  \label{fig:egitim-veri-matris}
\end{figure}

Üç soru, üç cevap:

\textbf{VTUAV çok büyük, maskelerini nerede kullanıyoruz?} VTUAV'ın 500
dizisinden yalnızca 100'ünde (VIS bölümü) gerçek instance mask var; kalanın
tek etiketi kare başına bir kutu. Maskeler aşama B'ye gidiyor, çünkü orada
hedef bir maskedir. Setin \emph{geri kalanı}, yani asıl hacim, aşama A'ya
gidiyor; orada hiçbir etiket okunmadığı için maskesiz olmak hiçbir şeye mal
olmuyor. Yani VTUAV iki aşamayı iki farklı yanıyla besliyor.

\textbf{SegFly ve Kust4K video değil, nasıl işe yarıyor?} Aşama A ve B'nin
ikisi de \textbf{tek kare} üzerinde çalışıyor. Video yalnızca aşama C'de
gerekiyor. İki set de hizalı RGB-T çifti taşıdığı için aşama A'yı besliyor,
dense mask taşıdığı için de aşama B'yi.

\textbf{Semantic set nasıl instance hedefi veriyor?} Aşama B'nin tamamı
bu: hedefi biz kuruyoruz. Thing sınıflarını ayırıp sınıf sınıf
connected component alıyoruz, her component bir prompt ve bir target
oluyor. SegFly hacmi getiriyor (Kust4K'nın dört katı kare, üç irtifa),
Kust4K en ince etiketi (dört thing sınıfı; SegFly'da yalnızca iki tane
var). Tek set kullanmak tek sensör, tek şehir ve tek etiketleme alışkanlığı
demek olurdu, o yüzden üçü tek havuzda birleşiyor ve bir batch üçünü aynı
adımda taşıyabiliyor.

Bir de \textbf{skorlama} ayrımı var. Semantic map'ten \emph{yeniden
kurulmuş} instance'ların ground truth'u kendisi belirsiz: ayrıştırma iki
arabayı kaynaştırdıysa ``doğru'' cevap tek blob olur ve onları doğru ayıran
model yanlış sayılır. O sayı modeli değil ayrıştırmayı ölçer. Bu yüzden
Kust4K ve SegFly yalnızca eğitime giriyor; skorlama VTUAV VIS'in çizilmiş
maskeleri üzerinden yapılıyor. Gürültülü bir target'tan öğrenmek normal,
gürültülü bir target'la not vermek değil.

Hizalı çift bir şeyi daha mümkün kılıyor. Güçlü bir RGB modeli RGB
yarısında maske üretip aynı maskeyi termal yarıya taşıyabilir; böylece
maskesiz kareler aşama B'ye veri olur, ve encoder iki modaliteyi birden
görür. Tek şartı hizalama, ve bu bir varsayım değil ölçülecek bir şey:
hedefler karenin yüzde biri kadar, birkaç piksellik sistematik bir kayma
maskeyi bozar. VTUAV'ın kendi belgesi iki modalitenin piksel piksel hizalı
\emph{olmadığını} söylüyor, o yüzden fikir setten sete ayrı ayrı
ölçülmeden uygulanmıyor.

VTUAV için iki ayar başka hiçbir yerde olmadığı kadar önemli, ikisi de
$1920\times1080$ video olmasından geliyor. \textbf{Kırpma oranı 512/1080:}
16:9 bir kareyi $512\times512$'ye sıkıştırmak aspect ratio'yu bozuyor ve
küçük bir aracı encoder'ın göremeyeceği boyutun altına indiriyor; kırpma
her iki yarıdan aynı normalize konumdaki kareyi alıyor, böylece hizalama
korunuyor. \textbf{Bir hücrelik tolerans:} VTUAV'ın iki modalitesi piksel
piksel hizalı değil (yazarların kendi notu), o yüzden her student pozisyonu
çevresindeki $3\times3$ komşuluğun en iyi teacher pozisyonuyla eşleşiyor.
Bedeli uzamsal kesinlik, o yüzden hizalı setlerde kapalı.

\bigskip
\noindent
Sıradaki adım bir kod değil, bir koşu. İki sayı bu bölümdeki her cümleyi ya
doğruluyor ya da çöpe atıyor: \textbf{kaynaşma oranı} (GPU'suz, dakikalar
içinde) ve \textbf{baseline'a karşı aşama A}. Aşama C, ikincisi gelmeden
başlamıyor.
